%% ma: L11-B16-0020
%% lop: 11
%% chuong: 5
%% bai: 16
%% dang: 11.16.5
%% loai: TN4
%% muc: TH
%% dap_an: "C"
%% nguon: "tai-lieu-hoc-tap-mon-toan-lop-11-theo-chuong-trinh-giao-duc-pho-thong-2018.pdf (Lớp Toán Thầy Duy, Phạm Lê Duy, THPT Chu Văn An), Chương 3 Bài 2, Phần 1, Câu 8 (tr.360)"
%% kien_thuc: "Lớp 11 Bài 16 (giới hạn của hàm số)"
%% trang_thai: chinh-thuc
%% ghi_chu: "Đã duyệt 03/11/2026."
\begin{ex}
Trong bốn giới hạn sau đây, giới hạn nào bằng $-\infty$?
\choice{$\lim\limits_{x\to+\infty}\dfrac{-3x+4}{x-2}$.}{$\lim\limits_{x\to2^-}\dfrac{-3x+4}{x-2}$.}{\True $\lim\limits_{x\to2^+}\dfrac{-3x+4}{x-2}$.}{$\lim\limits_{x\to-\infty}\dfrac{-3x+4}{x-2}$.}
\loigiai{\sloppy A, D bằng $-3$. B: tử $\to-2<0$, mẫu $\to0^-$ nên giới hạn $+\infty$. C: tử $\to-2<0$, mẫu $\to0^+$ nên giới hạn $-\infty$.}
\end{ex}
